% !PRINTED 418
\weberpage{418}
%<<CONT
\gesperrt{integral rational functions of $x$ with rational numerical coefficients.}

This important theorem was set up and first proved by \gesperrt{Gauss} for
the case that $n$ is a prime. Later many other proofs have been given, some
of them more general and simpler. We follow here a proof by
\gesperrt{Kronecker}.

In the function $X_n$, as we have seen above, the highest power of $x$ has
the coefficient $1$, while the remaining coefficients are integers. If we
now assume that $X_n$ can be decomposed into two factors with rational
coefficients

\weq{1}{X_n = \varphi(x)\,\psi(x)}

then the product of the coefficients of the highest powers of $x$ in
$\varphi$ and $\psi$ must be equal to $1$, and we may therefore assume that
these highest coefficients in $\varphi$ and $\psi$ themselves have the value
$1$ (since otherwise we would divide by the product of the two
coefficients). But then it follows from the theorem of \gesperrt{Gauss}
(\S.\,2) that the remaining coefficients of $\varphi(x)$ and $\psi(x)$ must
be \gesperrt{integers}, and hence that $\varphi(x)$ and $\psi(x)$, when an integer is
put for $x$, must themselves pass into integers.

If we therefore set $x = 1$, it follows that:
\[
\varphi(1)\,\psi(1) = X_n(1),
\]
and so by Theorem V. this is equal to $1$ or to a prime $p$. Accordingly,
of the two factors $\varphi(1)$, $\psi(1)$, one must certainly have the
value $\pm 1$, while the other is equal to $\pm 1$ or $\pm p$. We therefore
assume that

\weq{2}{\varphi(1) = \pm 1.}

The roots $r$ of the equation $X_n = 0$ now distribute themselves between
the two equations $\varphi = 0$, $\psi = 0$, and since neither of the two
factors $\varphi$, $\psi$ is independent of $x$, there must certainly be
one among the primitive $n$th roots of unity $r$, which we shall denote by
$\rho$, for which
\[
\varphi(\rho) = 0
\]
holds.

Now if $r$ denotes any primitive $n$th root of unity, then among the powers
\[
r,\ r^2,\ r^3\ \dots\ r^{n-1}
\]
all $n$th roots of unity with the exception of $1$, and therefore also
$\rho$, are contained; and it follows therefore that the product
%CONT>>
